\documentclass[11pt, a4paper]{article}

\usepackage[T1]{fontenc}
\usepackage[utf8]{inputenc}
\usepackage[a4paper, top=2.5cm, bottom=2.5cm, left=2.5cm, right=2.5cm]{geometry}
\usepackage{booktabs}
\usepackage{enumitem}
\usepackage[hidelinks]{hyperref}
\usepackage{parskip}
\usepackage{amsmath}
\usepackage{listings}
\usepackage{xcolor}

\definecolor{codigobg}{RGB}{245,245,245}
\definecolor{keyword}{RGB}{0,100,150}
\lstset{language=R, basicstyle=\ttfamily\small, backgroundcolor=\color{codigobg},
  keywordstyle=\color{keyword}\bfseries, frame=single, breaklines=true, showstringspaces=false}

\setlength{\parindent}{0pt}

\pagestyle{plain}

\begin{document}

\begin{center}
{\large\bfseries PONTIFICIA UNIVERSIDAD JAVERIANA}\\[3pt]
Facultad de Ciencias Econ\'omicas y Administrativas\\[2pt]
Estad\'istica 2026-II --- Prof.\ Jaime Polanco Jim\'enez\\[6pt]
{\LARGE\bfseries Problem Set 3}\\[3pt]
{\large Distribuci\'on Normal y Teorema Central del L\'imite}\\[6pt]
Asignado: Agosto 21 (Sesi\'on 3) \quad$|$\quad Entrega: Septiembre 11 (Sesi\'on 4)\\[3pt]
\textit{Plazo extendido: 3 semanas (Semana de Reflexi\'on + Parcial 1)}
\end{center}

\vspace{6pt}
\rule{\linewidth}{0.5pt}
\vspace{2pt}
{\centering\small Repositorio: \url{https://github.com/polanco-jaime/Curso-Estadistica-2026-3}\par}
\vspace{6pt}

\textbf{Instrucciones.} Este Problem Set es individual y se resuelve en R sobre
microdatos ICFES. La entrega es digital al inicio de la clase indicada.
El c\'odigo R debe incluirse.

\vspace{10pt}

\section*{Enunciado}

Este Problem Set tiene dos partes independientes: ejercicios con la distribuci\'on
normal y una simulaci\'on del Teorema Central del L\'imite (TCL).

\subsection*{Parte 1: Ejercicios con la Distribuci\'on Normal}

Usando la variable \texttt{MOD\_INGLES\_PUNT} (puntaje de ingl\'es en Saber Pro)
para estudiantes de Negocios Internacionales (NI):

\begin{enumerate}[label=\alph*)]
  \item \textbf{Par\'ametros:} calcule la media $\mu$ y desviaci\'on est\'andar
        $\sigma$ de \texttt{MOD\_INGLES\_PUNT}. Asuma que esta variable se distribuye
        aproximadamente normal.

  \item \textbf{C\'alculo de probabilidades con \texttt{pnorm()}:}
        \begin{itemize}
          \item ¿Qu\'e proporci\'on de estudiantes obtiene un puntaje mayor a 150?
          \item ¿Qu\'e proporci\'on obtiene un puntaje entre 120 y 160?
          \item ¿Cu\'al es la probabilidad de obtener un puntaje menor a 100?
        \end{itemize}

  \item \textbf{C\'alculo de cuantiles con \texttt{qnorm()}:}
        \begin{itemize}
          \item ¿Qu\'e puntaje deja por debajo al 25\% de los estudiantes?
          \item ¿Qu\'e puntaje deja por debajo al 90\% de los estudiantes?
          \item Calcule el rango intercuart\'ilico (IQR) usando cuantiles.
        \end{itemize}

  \item \textbf{Visualizaci\'on:} grafique el histograma de \texttt{MOD\_INGLES\_PUNT}
        y superponga la curva de densidad normal te\'orica $N(\mu, \sigma^2)$.
        Comente sobre el ajuste.
\end{enumerate}

\subsection*{Parte 2: Simulaci\'on del Teorema Central del L\'imite}

Utilizando los microdatos de Saber 11, realice una simulaci\'on del TCL con la
variable \texttt{PUNT\_INGLES} (puntaje de ingl\'es).

\begin{enumerate}[label=\alph*), resume]
  \setcounter{enumi}{4}
  \item \textbf{Poblaci\'on:} describa la distribuci\'on poblacional de
        \texttt{PUNT\_INGLES} en Saber 11. Calcule $\mu$ y $\sigma$ poblacionales.
        Grafique la distribuci\'on con un histograma.

  \item \textbf{Simulaci\'on del TCL:} para cada tama\~no de muestra
        $n \in \{30,\, 100,\, 500\}$:
        \begin{itemize}
          \item Extraiga 2000 muestras aleatorias de tama\~no $n$ de la poblaci\'on.
          \item Calcule la media muestral $\bar{X}$ de cada muestra.
          \item Grafique la distribuci\'on de las 2000 medias muestrales.
        \end{itemize}

  \item \textbf{Convergencia a la normalidad:} para cada $n$, superponga sobre el
        histograma de medias muestrales la curva te\'orica
        $N\!\left(\mu,\, \sigma^2/n\right)$.
        Comente sobre la convergencia a medida que $n$ crece.

  \item \textbf{Error est\'andar:} para cada $n$, calcule:
        \begin{itemize}
          \item El error est\'andar te\'orico: $\text{SE}_{\text{te\'or}} = \sigma / \sqrt{n}$
          \item El error est\'andar simulado: la desviaci\'on est\'andar de las 2000 medias
                muestrales.
        \end{itemize}
        Compare ambos valores en una tabla. ¿Qu\'e tan cerca est\'an?

  \item \textbf{Conclusiones:} explique en sus propias palabras qu\'e demuestra
        esta simulaci\'on sobre el TCL y por qu\'e es relevante para la inferencia
        estad\'istica.
\end{enumerate}

\vspace{10pt}
\rule{\linewidth}{0.4pt}
\begin{center}
{\small Cada entrega completa suma $+0.0625$ a la nota final.}
\end{center}

\end{document}
